\documentclass[11pt,a4paper]{article}
\usepackage[margin=18mm,top=22mm,bottom=20mm]{geometry}
\usepackage{fontspec}\setmainfont{Latin Modern Roman}
\usepackage{amsmath,amssymb,tikz,graphicx,fancyhdr,array,needspace}
\usetikzlibrary{patterns,arrows.meta,decorations.pathmorphing,calc}
\setlength{\headheight}{14pt}
\usetikzlibrary{arrows.meta}
\pagestyle{fancy}\fancyhf{}\renewcommand{\headrulewidth}{0.4pt}
\fancyhead[L]{\textbf{M-03}}\fancyhead[C]{\itshape Fizika (OliyFizika.uz mock test)}\fancyhead[R]{\textbf{Mock №3}}
\fancyfoot[C]{\thepage}
\setlength{\parindent}{0pt}\setlength{\parskip}{3pt}
\newcommand{\opt}[2]{\hspace*{4mm}#1)\ #2\par}
\newcommand{\blank}{\rule{5cm}{0.4pt}}
\begin{document}

\begin{center}{\Large\textbf{FIZIKA}}\end{center}
\noindent\begin{minipage}{\linewidth}
\textbf{1.}\ Moddiy nuqta X öqi böylab musbat yönalişda 10 m/s boşlanğiç tezlik bilan tekis özgaruvçan harakatlanmoqda. Moddiy nuqta qanday maksimal tezlaniş (m/s$^{2}$) bilan harakatlanganda, boşlanğiç 2 s içida X öqining 10 m uzunlikdagi kesmasida mavjud böladi?\par\smallskip
\opt{A}{2}
\opt{B}{4}
\opt{C}{-5}
\opt{D}{5}
\end{minipage}\vspace{9pt}\par
\noindent\begin{minipage}{\linewidth}
\textbf{2.}\ Yer sirtidagi jism R/10 balandlikka kötarilganda bajarilgan işni toping. R-Yer radiusi, g-Yer sirtidagi erkin tuşiş tezlanişi.\par\smallskip
\opt{A}{mgR/10}
\opt{B}{10mgR}
\opt{C}{mgR/11}
\opt{D}{mgR/9}
\end{minipage}\vspace{9pt}\par
\noindent\begin{minipage}{\linewidth}
\textbf{3.}\ Gorizontal silliq stol ustida qoğoz yotibdi. Qoğoz ustiga esa suv solingan stakan qöyilgan. Qoğozni stakan tagidan qanday eng kiçik tezlaniş bilan tortib olsak, stakan joyida qoladi (m/s$^{2}$). Qoğoz va stakan orasidagi işqalaniş koeffitsiyenti 0,3 ga teng. g=10 N/kg.\par\smallskip
\opt{A}{5}
\opt{B}{3}
\opt{C}{4}
\opt{D}{6}
\end{minipage}\vspace{9pt}\par
\noindent\begin{minipage}{\linewidth}
\textbf{4.}\ Bekatdan qözğalayotgan poyezdning birinçi vagoni poyezd joyidan qözğalgunga qadar şu vagonning boşida turgan kuzatuvçining yonidan 3 s da ötadi. 9 ta vagondan iborat butun poyezd qança vaqtda kuzatuvçining yonidan ötadi (s)? Vagonlar orasidagi oraliqlarni hisobga olmang.\par\smallskip
\opt{A}{10}
\opt{B}{9}
\opt{C}{3}
\opt{D}{27}
\end{minipage}\vspace{9pt}\par
\noindent\begin{minipage}{\linewidth}
\textbf{5.}\ Moddiy nuqta $x = A\cos(\pi t/6 + \pi/2)$ qonun böyiça tebranmoqda. Vaqt ölçovi $t$ = 0 dan boşlansa qanday $t$ vaqtda (s) moddiy nuqta bosgan yöl 0,5$A$ ga yetadi?\par\smallskip
\opt{A}{2}
\opt{B}{1}
\opt{C}{3}
\opt{D}{4}
\end{minipage}\vspace{9pt}\par
\noindent\begin{minipage}{\linewidth}
\textbf{6.}\ Erkin qöyib yuborilgan jism 5 s da yerga tuşdi. U qanday balandlikdan (m) tuşganini aniqlang. g=10 m/s$^{2}$\par\smallskip
\opt{A}{180}
\opt{B}{80}
\opt{C}{45}
\opt{D}{125}
\end{minipage}\vspace{9pt}\par
\noindent\begin{minipage}{\linewidth}
\textbf{7.}\ Tik yuqoriga 10 m/s tezlik bilan F otilgan 1 kg massali jism yerga qaytib tuşdi. Oğirlik kuçi maydonida bajarilgan işni hisoblang (J).\par\smallskip
\opt{A}{50}
\opt{B}{100}
\opt{C}{0}
\opt{D}{75}
\end{minipage}\vspace{9pt}\par
\noindent\begin{minipage}{\linewidth}
\textbf{8.}\ Baland qoyadan 0,1 s vaqt intervali bilan tomçi tommoqda. 10-tomçi tomgandan söng 2- va 5-tomçilar orasidagi masofa qanday böladi (m)? g=10 m/s$^{2}$.\par\smallskip
\opt{A}{1,95}
\opt{B}{2,95}
\opt{C}{3}
\opt{D}{2,85}
\end{minipage}\vspace{9pt}\par
\noindent\begin{minipage}{\linewidth}
\textbf{9.}\ Vaznsiz blok 40 N kuç bilan yuqoriga kötarilmoqda. Blok orqali ötkazilgan ipning bir uçiga 1 kg, ikkinçi uçiga 2 kg massali yuklar osilgan. Birinçi yukning yerga nisbatan tezlanişi (m/s$^{2}$) topilsin. $g$=10 m/s$^{2}$\par\smallskip
\opt{A}{10}
\opt{B}{20}
\opt{C}{14}
\opt{D}{10/3}
\end{minipage}\vspace{9pt}\par
\noindent\begin{minipage}{\linewidth}
\textbf{10.}\ Yer sirtidan 0,1R balandlikda turgan jism Yerni butunlay tark etişi uçun qanday minimal kinetik energiyaga ega bölişi kerak? R-Yer radiusi, g-Yer sirtidagi erkin tuşiş tezlanişi.\par\smallskip
\opt{A}{mgR}
\opt{B}{mgR/2}
\opt{C}{2mgR}
\opt{D}{mgR/1,1}
\end{minipage}\vspace{9pt}\par
\noindent\begin{minipage}{\linewidth}
\textbf{11.}\ Hajmi V = 3 m$^{3}$ bölgan ballonga N = 0,6$\cdot10^{18}$ ta azot molekulalari qamalgan bölsa, undagi molekulalar konsentratsiyasini (m$-$3) toping.\par\smallskip
\opt{A}{5$\cdot10^{18}$}
\opt{B}{0,2$\cdot10^{18}$}
\opt{C}{1,8$\cdot10^{18}$}
\opt{D}{2$\cdot10^{18}$}
\end{minipage}\vspace{9pt}\par
\noindent\begin{minipage}{\linewidth}
\textbf{12.}\ Hajmi 3 l, temperaturasi 90 $^\circ$C bölgan suvga hajmi 5 l, temperaturasi esa 10 $^\circ$C bölgan suv aralaştirildi. Natijaviy temperaturani ($^\circ$C) aniqlang.\par\smallskip
\opt{A}{120}
\opt{B}{50}
\opt{C}{40}
\opt{D}{60}
\end{minipage}\vspace{9pt}\par
\noindent\begin{minipage}{\linewidth}
\textbf{13.}\ Havo molekulalari Yer atmosferasini tark etişlari uçun Yerda harorat kamida qanday bölişi lozim? Havoning molyar massasi $M$, Yer uçun ikkinçi kosmik tezlik $u$.\par\smallskip
\opt{A}{$3Mu^2/2R$}
\opt{B}{$Mu^2/3R$}
\opt{C}{$Mu/3R$}
\opt{D}{$Mu^2/2R$}
\end{minipage}\vspace{9pt}\par
\noindent\begin{minipage}{\linewidth}
\textbf{14.}\ Normal şaroitda har qanday gazning birlik hajmdagi molekulalar soni nimaga teng?\par\smallskip
\opt{A}{Avogadro soniga}
\opt{B}{Faradey soniga}
\opt{C}{Loşmidt soniga}
\opt{D}{Boltsman soniga}
\end{minipage}\vspace{9pt}\par
\noindent\begin{minipage}{\linewidth}
\textbf{15.}\ Massasi 16 g bölgan geliy gazining hajmi $p\sqrt[3]{V}=\mathrm{const}$ qonuniyat böyiça 8 marta ortdi. Uning boşlanğiç temperaturasi 150 K bölsa, içki energiyasi qanday özgargan?\par\smallskip
\opt{A}{12,5 kJ ga kamaygan}
\opt{B}{22,5 kJ ga kamaygan}
\opt{C}{12,5 kJ ga ortgan}
\opt{D}{22,5 kJ ga ortgan}
\end{minipage}\vspace{9pt}\par
\noindent\begin{minipage}{\linewidth}
\textbf{16.}\ $C$ siğimli birinçi kondensatordagi kuçlaniş $U$, $2C$ siğimli ikkinçi kondensatordagi kuçlaniş $2U$. Kondensatorlarning qoplamalari töğri (zaryadlari moslab) ulansa, batareyaning natijaviy energiyasi qanday böladi?\par\smallskip
\opt{A}{$32CU^2/5$}
\opt{B}{$16CU^2/5$}
\opt{C}{$20CU^2/3$}
\opt{D}{$25CU^2/6$}
\end{minipage}\vspace{9pt}\par
\noindent\begin{minipage}{\linewidth}
\textbf{17.}\ $6{,}28\cdot10^{-15}$ C/m$^{2}$ sirtiy ziçlikli metall şarning zaryadi $8\cdot10^{-10}$ C ga teng. Şarning radiusini toping (m). $\pi^2=10$.\par\smallskip
\opt{A}{100}
\opt{B}{10}
\opt{C}{1}
\opt{D}{1000}
\end{minipage}\vspace{9pt}\par
\noindent\begin{minipage}{\linewidth}
\textbf{18.}\ FIK i 90 \% bölgan elektronasosdan 200 V kuçlanişda 10 A tok ötmoqda. Uşbu nasos 1 soatda qanday hajmdagi suvni 36 m balandlikka kötaradi (m$^{3}$)? $\rho$=1000 kg/m$^{3}$, g=10 N/kg.\par\smallskip
\opt{A}{27}
\opt{B}{9}
\opt{C}{18}
\opt{D}{36}
\end{minipage}\vspace{9pt}\par
\noindent\begin{minipage}{\linewidth}
\textbf{19.}\ Qoplamalari orasidagi masofa 1,5 mm bölgan kondensator 100 V kuçlanişga ulangan. Qoplamalardagi zaryad sirt ziçligini (nC/m$^{2}$) aniqlang.\par\smallskip
\opt{A}{230}
\opt{B}{540}
\opt{C}{590}
\opt{D}{420}
\end{minipage}\vspace{9pt}\par
\noindent\begin{minipage}{\linewidth}
\textbf{20.}\ 2 kondensatorni ketma-ket ulab batareya hosil qilingan. Batareya akkumulyatorga ulanib zaryadlangan. Batareya $W$ energiya töplagan bölsa, akkumulyator qança energiya sarflagan?\par\smallskip
\opt{A}{$3W$}
\opt{B}{$1{,}5W$}
\opt{C}{$2W$}
\opt{D}{$W$}
\end{minipage}\vspace{9pt}\par
\noindent\begin{minipage}{\linewidth}
\textbf{21.}\ Quyidagi fikrlarning qaysilari töğri? 1) vakuumli diodning zaryad taşuvçilari elektronlar, musbat va manfiy ionlardan iborat; 2) $p$ turdagi yarimötkazgiçlarda asosiy zaryad taşuvçilar-elektronlar; 3) töğri $p-n$ ötişda yarim ötkazgiçli diodning berkituvçi qatlami torayadi; 4) yarimötkazgiçlarda elektron va kovak uçraşganda neytral atom hosil böladi\par\smallskip
\opt{A}{1, 4}
\opt{B}{1, 2}
\opt{C}{2, 3}
\opt{D}{3, 4}
\end{minipage}\vspace{9pt}\par
\noindent\begin{minipage}{\linewidth}
\textbf{22.}\ Yukçasi diamagnit moddadan bölgan prujinali mayatnik vertikal tebranişlarni amalga oşirmoqda. Agar mayatnik tagiga magnit joylaştirilsa, uning tebraniş davri qanday özgaradi? Magnitning prujina bilan özaro ta’sirini hisobga olmang.\par\smallskip
\opt{A}{ortadi}
\opt{B}{kamayadi}
\opt{C}{özgarmaydi}
\opt{D}{avval ortib, keyin kamayadi}
\end{minipage}\vspace{9pt}\par
\noindent\begin{minipage}{\linewidth}
\textbf{23.}\ Orasi l=1 m bölgan gorizontal relslarda m=200 g massali ötkazgiç sterjen yotibdi (sterjen relslarga perpendikulyar). Butun sistema induksiyasi B=0,5 T bölgan perpendikulyar magnit maydonda joylaşgan. Agar sterjen orqali I=4 A tok ötkazilsa, u relslarda Amper kuçi ta’sirida a=6 m/s$^{2}$ tezlaniş bilan harakatlanadi. Sterjenning relslarga sirpaniş işqalaniş koeffitsenti $\mu$ ni toping. g=10 N/kg.\par\smallskip
\opt{A}{0,4}
\opt{B}{0,5}
\opt{C}{0,1}
\opt{D}{0,3}
\end{minipage}\vspace{9pt}\par
\noindent\begin{minipage}{\linewidth}
\textbf{24.}\ Qizuvçi simlari volframdan yasalgan ikki lampa özgarmas kuçlaniş manbayiga parallel ulanganida birinçi lampa ravşanroq yondi. Agar ular tarmoqqa ketma-ket ulansa, qanday yonadi?\par\smallskip
\opt{A}{ikkisi birdek yonadi}
\opt{B}{ikkinçisi ravşanroq yonadi}
\opt{C}{birinçisi ravşanroq yonadi}
\opt{D}{avval ikkinçisi, söngra birinçisi ravşanroq yonadi}
\end{minipage}\vspace{9pt}\par
\noindent\begin{minipage}{\linewidth}
\textbf{25.}\ $q_1=-40$ va $q_2=+75$ nC nuqtaviy zaryadlar bir-biridan 80 cm masofada joylaşgan. Birinçi zaryaddan $r_1$=40 cm va ikkinçi zaryaddan $r_2$=50 cm masofada turgan nuqtadagi maydon potensiali $\varphi$ (V) topilsin.\par\smallskip
\opt{A}{450}
\opt{B}{500}
\opt{C}{600}
\opt{D}{350}
\end{minipage}\vspace{9pt}\par
\noindent\begin{minipage}{\linewidth}
\textbf{26.}\ Tölqin uzunligi 500 nm bölgan monoxromatik yoruğlik tölqini davri 2 $\mu$m bölgan difraksion panjaraga tik tuşmoqda. Ikkinçi tartibli difraksion maksimum qanday burçak ostida körinadi?\par\smallskip
\opt{A}{300}
\opt{B}{600}
\opt{C}{450}
\opt{D}{150}
\end{minipage}\vspace{9pt}\par
\noindent\begin{minipage}{\linewidth}
\textbf{27.}\ Suyuqlik sirtiga yoruğlik nuri 36$^\circ$ burçak ostida tuşmoqda. Agar siniş burçagi 24$^\circ$ bölsa, singan va qaytgan nurlar orasidagi burçakni toping.\par\smallskip
\opt{A}{120$^\circ$}
\opt{B}{60$^\circ$}
\opt{C}{54$^\circ$}
\opt{D}{72$^\circ$}
\end{minipage}\vspace{9pt}\par
\noindent\begin{minipage}{\linewidth}
\textbf{28.}\ Yassi közguga nur 30$^\circ$ burçak ostida tuşmoqda. Tuşgan va qaytgan nurlar orasidagi burçakni aniqlang.\par\smallskip
\opt{A}{30$^\circ$}
\opt{B}{45$^\circ$}
\opt{C}{60$^\circ$}
\opt{D}{90$^\circ$}
\end{minipage}\vspace{9pt}\par
\noindent\begin{minipage}{\linewidth}
\textbf{29.}\ Vodorod atomi elektronining ikkinçi orbitadagi energiyasi E ga teng bölsa, törtinçi orbitadagi energiyasini toping.\par\smallskip
\opt{A}{E/2}
\opt{B}{E/8}
\opt{C}{E/4}
\opt{D}{E/3}
\end{minipage}\vspace{9pt}\par
\noindent\begin{minipage}{\linewidth}
\textbf{30.}\ Ikkita vodorod atomi berilgan, ularning energiyalari quyidagiça boğlangan: $E_1 = 9E_2$. Ulardagi elektronlarning impuls momentlari $M_1$ va $M_2$ qanday munosabatda böladi?\par\smallskip
\opt{A}{$M_1 = 3M_2$}
\opt{B}{$M_1 = 9M_2$}
\opt{C}{$9M_1 = M_2$}
\opt{D}{$3M_1 = M_2$}
\end{minipage}\vspace{9pt}\par
\noindent\begin{minipage}{\linewidth}
\textbf{31.}\ Yadroda neytron protonga aylandi. Bunda antineytrino bilan yana qaysi zarra ajralib çiqadi?\par\smallskip
\opt{A}{alfa zarra}
\opt{B}{elektron}
\opt{C}{pozitron}
\opt{D}{gamma kvant}
\end{minipage}\vspace{9pt}\par
\noindent\begin{minipage}{\linewidth}
\textbf{32.}\ Yer sirtidagi jism R/n balandlikka kötarilganda bajarilgan işni toping. R-Yer radiusi, g-Yer sirtidagi erkin tuşiş tezlanişi.\par\smallskip
\opt{A}{mgR/(n-1)}
\opt{B}{nmgR}
\opt{C}{mgR/n}
\opt{D}{mgR/(n+1)}
\end{minipage}\vspace{9pt}\par
\newpage\noindent\textbf{33–35-test topşiriqlariga mos keluvçi javoblarni (A–F) javob variantlaridan tanlang.}\par\smallskip
Rasmda EYuK manbai va kondensatorlardan iborat zanjir tasvirlangan: $\varepsilon=12$ V, $C_1=6\ \mu$F, $C_2=3\ \mu$F, $C_3=3\ \mu$F.\par\medskip
\begin{center}\begin{tikzpicture}[x=1cm,y=1cm,line width=0.5pt,>=Stealth,font=\small]
\draw (0,0)--(0,1.0) (0,1.3)--(0,2.4)--(1.5,2.4); \draw (-0.3,1.3)--(0.3,1.3); \draw (-0.15,1.0)--(0.15,1.0); \node[left] at (-0.35,1.15) {$\varepsilon$};
\draw (1.5,2.1)--(1.5,2.7) (1.8,2.1)--(1.8,2.7); \draw (1.8,2.4)--(5.0,2.4); \node[above] at (1.65,2.75) {$C_1$};
\draw (0,0)--(5.0,0);
\foreach \x/\n in {3.3/C_2,5.0/C_3}{\draw (\x,2.4)--(\x,1.35) (\x,1.05)--(\x,0); \draw (\x-0.3,1.35)--(\x+0.3,1.35) (\x-0.3,1.05)--(\x+0.3,1.05); \node[right] at (\x+0.3,1.2) {$\n$};}
\fill (3.3,2.4) circle(1.2pt) (3.3,0) circle(1.2pt);
\end{tikzpicture}\end{center}
\noindent\fbox{\begin{minipage}{0.64\linewidth}\vspace{4pt}\textbf{33.}\ Manbadan zanjirga ötgan zaryad miqdorini ($\mu$C) toping.\\[6pt]\textbf{34.}\ $C_1$ kondensatordagi kuçlanişni (V) toping.\\[6pt]\textbf{35.}\ Zanjirning töliq energiyasini ($\mu$J) toping.\vspace{4pt}\end{minipage}\vrule\hspace{6pt}\begin{minipage}{0.28\linewidth}\vspace{4pt}A)\ 6\\B)\ 18\\C)\ 36\\D)\ 72\\E)\ 108\\F)\ 216\vspace{4pt}\end{minipage}}\par\vspace{14pt}
\Needspace{9\baselineskip}
\noindent\textbf{36.}\ Jismning havodagi oğirligi 10 N, suvga töliq botirilgandagi oğirligi 6 N. $\rho_{suv}=1000$ kg/m$^3$, $g=10$ m/s$^2$.\par\smallskip
\medskip\noindent\textbf{a)\ Jismning hajmini (cm$^3$) toping.}\par
\noindent\textit{Javob}: a)\ \blank\par
\medskip\noindent\textbf{b)\ Jism moddasining ziçligini (kg/m$^3$) toping.}\par
\noindent\textit{Javob}: b)\ \blank\par
\medskip\noindent\textbf{Diqqat!} Javoblaringizni javoblar varaqasiga köçirib yozing.\par\vspace{10pt}
\Needspace{9\baselineskip}
\noindent\textbf{37.}\ Jism gorizontga $37^\circ$ burçak ostida 50 m/s tezlik bilan otildi. $\sin37^\circ=0{,}6$, $\cos37^\circ=0{,}8$, $g=10$ m/s$^2$.\par\smallskip
\medskip\noindent\textbf{a)\ Jismning eng katta kötariliş balandligini (m) toping.}\par
\noindent\textit{Javob}: a)\ \blank\par
\medskip\noindent\textbf{b)\ Jismning uçiş uzoqligini (m) toping.}\par
\noindent\textit{Javob}: b)\ \blank\par
\medskip\noindent\textbf{Diqqat!} Javoblaringizni javoblar varaqasiga köçirib yozing.\par\vspace{10pt}
\Needspace{9\baselineskip}
\noindent\textbf{38.}\ Bikrligi 800 N/m bölgan prujinali töpponça prujinasi 5 cm siqilgan. Töpponçadan 20 g massali şarça tik yuqoriga otildi. Havoning qarşiligi hisobga olinmasin, $g=10$ m/s$^2$.\par\smallskip
\medskip\noindent\textbf{a)\ Şarçaning otiliş tezligini (m/s) toping.}\par
\noindent\textit{Javob}: a)\ \blank\par
\medskip\noindent\textbf{b)\ Şarçaning kötariliş balandligini (m) toping.}\par
\noindent\textit{Javob}: b)\ \blank\par
\medskip\noindent\textbf{Diqqat!} Javoblaringizni javoblar varaqasiga köçirib yozing.\par\vspace{10pt}
\Needspace{9\baselineskip}
\noindent\textbf{39.}\ 100 kPa bosimdagi 6 L gaz izotermik ravişda 2 L gaça siqildi, söngra özgarmas hajmda 27 $^\circ$C dan 127 $^\circ$C gaça qizdirildi.\par\smallskip
\medskip\noindent\textbf{a)\ Izotermik siqilişdan keyingi bosimni (kPa) toping.}\par
\noindent\textit{Javob}: a)\ \blank\par
\medskip\noindent\textbf{b)\ Qizdirişdan keyingi gaz bosimini (kPa) toping.}\par
\noindent\textit{Javob}: b)\ \blank\par
\medskip\noindent\textbf{Diqqat!} Javoblaringizni javoblar varaqasiga köçirib yozing.\par\vspace{10pt}
\Needspace{9\baselineskip}
\noindent\textbf{40.}\ Quvvati 2 kW bölgan elektr çoynakda 1,5 L suv 20 $^\circ$C dan 100 $^\circ$C gaça isitildi. Çoynakning FIK 84\%. $c_{suv}=4200$ J/(kg$\cdot$K).\par\smallskip
\medskip\noindent\textbf{a)\ Suvni isitiş uçun zarur bölgan issiqlik miqdorini (kJ) toping.}\par
\noindent\textit{Javob}: a)\ \blank\par
\medskip\noindent\textbf{b)\ Suvni isitiş uçun ketgan vaqtni (s) toping.}\par
\noindent\textit{Javob}: b)\ \blank\par
\medskip\noindent\textbf{Diqqat!} Javoblaringizni javoblar varaqasiga köçirib yozing.\par\vspace{10pt}
\Needspace{9\baselineskip}
\noindent\textbf{41.}\ 200 öramli ğaltakning köndalang kesim yuzasi 50 cm$^2$. Ğaltak öqi böylab yönalgan magnit maydon induksiyasi 0,02 s davomida 0,1 T dan 0,5 T gaça bir tekis ortdi. Ğaltak qarşiligi 4 $\Omega$, uçlari tutaştirilgan.\par\smallskip
\medskip\noindent\textbf{a)\ Ğaltakda hosil bölgan induksiya EYuKni (V) toping.}\par
\noindent\textit{Javob}: a)\ \blank\par
\medskip\noindent\textbf{b)\ Ğaltakdagi induksion tok kuçini (A) toping.}\par
\noindent\textit{Javob}: b)\ \blank\par
\medskip\noindent\textbf{Diqqat!} Javoblaringizni javoblar varaqasiga köçirib yozing.\par\vspace{10pt}
\Needspace{9\baselineskip}
\noindent\textbf{42.}\ Massasi 2 g bölgan $5\ \mu$C zaryadli şarça ingiçka ipga osilgan va kuçlanganligi 3000 V/m bölgan gorizontal bir jinsli elektr maydonda muvozanatda turibdi. $g=10$ m/s$^2$.\par\smallskip
\medskip\noindent\textbf{a)\ Şarçaga ta’sir qiluvçi elektr kuçini (mN) toping.}\par
\noindent\textit{Javob}: a)\ \blank\par
\medskip\noindent\textbf{b)\ Ipning taranglik kuçini (mN) toping.}\par
\noindent\textit{Javob}: b)\ \blank\par
\medskip\noindent\textbf{Diqqat!} Javoblaringizni javoblar varaqasiga köçirib yozing.\par\vspace{10pt}
\Needspace{9\baselineskip}
\noindent\textbf{43.}\ Şişaning sindiriş körsatkiçi 2 ga teng. $c=3\cdot10^8$ m/s.\par\smallskip
\medskip\noindent\textbf{a)\ Şişa–havo çegarasida töliq içki qaytişning çegaraviy burçagini (gradusda) toping.}\par
\noindent\textit{Javob}: a)\ \blank\par
\medskip\noindent\textbf{b)\ Yoruğlikning şişadagi tezligini ($10^8$ m/s) toping.}\par
\noindent\textit{Javob}: b)\ \blank\par
\medskip\noindent\textbf{Diqqat!} Javoblaringizni javoblar varaqasiga köçirib yozing.\par\vspace{10pt}
\Needspace{9\baselineskip}
\noindent\textbf{44.}\ Fokus masofasi 20 cm bölgan soçuvçi linzadan 20 cm masofada buyum turibdi.\par\smallskip
\medskip\noindent\textbf{a)\ Linzadan tasvirgaça masofani (cm) toping.}\par
\noindent\textit{Javob}: a)\ \blank\par
\medskip\noindent\textbf{b)\ Linzaning çiziqli kattalaştirişini toping.}\par
\noindent\textit{Javob}: b)\ \blank\par
\medskip\noindent\textbf{Diqqat!} Javoblaringizni javoblar varaqasiga köçirib yozing.\par\vspace{10pt}
\Needspace{9\baselineskip}
\noindent\textbf{45.}\ Geliy-4 yadrosining massasi 4,00151 m.a.b., proton massasi 1,00728 m.a.b., neytron massasi 1,00866 m.a.b. 1 m.a.b. 931,5 MeV energiyaga mos keladi.\par\smallskip
\medskip\noindent\textbf{a)\ Yadroning boğlaniş energiyasini (MeV) toping.}\par
\noindent\textit{Javob}: a)\ \blank\par
\medskip\noindent\textbf{b)\ Bitta nuklonga töğri keladigan soliştirma boğlaniş energiyasini (MeV) toping.}\par
\noindent\textit{Javob}: b)\ \blank\par
\medskip\noindent\textbf{Diqqat!} Javoblaringizni javoblar varaqasiga köçirib yozing.\par\vspace{10pt}
\end{document}